% =====================================================================
%  Notes PDF preamble — shared by every Quarto book in this repository.
%
%  Aesthetic: clean, minimalist, low-chrome. A single accent colour and a
%  thin left rule are the only ornamental elements. No boxes, no shading,
%  no headers/footers louder than the body.
%
%  Provides:
%   * Geometry, fonts, microtypography.
%   * Theorem-like environments (theorem, proposition, lemma, corollary,
%     definition, example, exercise, solution, remark, note) with a
%     unified left-rule style via thmtools + mdframed.
%   * `codeoutput` environment for runtime output (left rule, small mono).
%   * Math display tweaks for matrices / vectors / aligned blocks
%     (cleaner spacing — material for the next iteration).
% =====================================================================

% --- packages ---------------------------------------------------------
% NOTE: geometry, hyperref and xcolor are loaded by Quarto itself via the
% YAML front matter; reloading them here triggers "Option clash" errors.
\usepackage{microtype}
\usepackage{amsmath,amssymb,amsthm,mathtools}
\usepackage{siunitx}   % \num{…}, \SI{…}{…}, \si{…} — physics/data notes use these
\usepackage{cancel}    % \cancel{…} for striking through algebraic terms
% `physics` was here but conflicts with \mathcal{X}(y) inside \(...\).
% Provide the macros we actually use with plain LaTeX equivalents.
\providecommand{\abs}[1]{\left\lvert #1 \right\rvert}
\providecommand{\norm}[1]{\left\lVert #1 \right\rVert}
\providecommand{\vb}[1]{\boldsymbol{#1}}
\providecommand{\vec}[1]{\boldsymbol{#1}}
\providecommand{\dv}[2]{\frac{d#1}{d#2}}
\providecommand{\pdv}[2]{\frac{\partial #1}{\partial #2}}
\usepackage{dsfont}    % \mathds{1} indicator function used in ML notes
\usepackage{algorithm} % algorithm float used in ML notes
\usepackage{algpseudocode}
\usepackage{pdfpages}  % \includepdf{…} — used by 1sem/MAPDA/notes/chapters/cheatsheet.qmd
                       % to append the standalone landscape cheatsheet at the
                       % end of the downloadable PDF.

% ── TikZ + CircuitiKZ for diagrams (logic gates, FSMs, K-maps, etc.) ──
\usepackage{tikz}
\usepackage{circuitikz}
\usetikzlibrary{calc, arrows.meta, positioning, shapes.geometric,
                shapes.misc, fit, decorations.pathreplacing,
                decorations.pathmorphing, automata, matrix, chains}
\ctikzset{logic ports=ieee, logic ports/scale=0.85}

% ── Reusable TikZ styles for boolean / circuit / FSM diagrams ────────
\tikzset{
  > = stealth,
  thick wire/.style   = {thick, line cap=rect},
  state/.style        = {draw, circle, minimum size=8mm, font=\small},
  initial/.style      = {state, double},
  accepting/.style    = {state, fill=ink!8},
  kmap cell/.style    = {draw=rule, minimum size=8mm, font=\small},
  kmap label/.style   = {font=\scriptsize, inner sep=2pt},
  endpoint/.style     = {-{Circle},shorten >=-2pt},
}

% Karnaugh map skeleton (2 vars):  \kmapTwo{a}{b}{c}{d}
\providecommand{\kmapTwo}[4]{%
  \begin{tikzpicture}[baseline=(c.base)]
    \node[kmap cell] (c00) at (0,1) {#1};
    \node[kmap cell] (c01) at (1,1) {#2};
    \node[kmap cell] (c10) at (0,0) {#3};
    \node[kmap cell] (c11) at (1,0) {#4};
    \node[kmap label, above=1pt of c01] {$B$};
    \node[kmap label, left =1pt of c00] {$A$};
    \node (c) at (0.5,0.5) {};
  \end{tikzpicture}}

% Author-defined colour macros found in MAPDA / MTP source: \red{…},
% \green{…}, \orange{…}, \blue{…}. Map them onto xcolor.
\providecommand{\red}   [1]{\textcolor{red}{#1}}
\providecommand{\green} [1]{\textcolor{green!60!black}{#1}}
\providecommand{\orange}[1]{\textcolor{orange}{#1}}
\providecommand{\blue}  [1]{\textcolor{blue!70!black}{#1}}
\providecommand{\purple}[1]{\textcolor{purple}{#1}}
\providecommand{\boldred}   [1]{\textcolor{red}    {\boldsymbol{#1}}}
\providecommand{\boldgreen} [1]{\textcolor{green!60!black}{\boldsymbol{#1}}}
\providecommand{\boldorange}[1]{\textcolor{orange} {\boldsymbol{#1}}}
\providecommand{\boldblue}  [1]{\textcolor{blue!70!black}{\boldsymbol{#1}}}
\providecommand{\boldpurple}[1]{\textcolor{purple} {\boldsymbol{#1}}}

% ── Other helpers from MAPDA/MAPDB's commands.tex ────────────────────
% The originals use a few tiny conveniences that we recreate in plain
% LaTeX so chapters compile without pulling in their custom package set.
\providecommand{\q}[1]{``#1''}                  % quoted text
\providecommand{\sprod}[2]{\left\langle \boldsymbol{#1},\boldsymbol{#2}\right\rangle}
\providecommand{\mean}[1]{\left\langle #1 \right\rangle}
\DeclareMathOperator*{\argmax}{arg\,max}
\DeclareMathOperator*{\argmin}{arg\,min}
\DeclareMathOperator*{\mslim}{ms\text{-}lim}

% Table helpers used here and there in the originals.
\providecommand{\STAB}[1]{\begin{tabular}{@{}c@{}}#1\end{tabular}}
\providecommand{\thickhline}{\noalign{\hrule height1.2pt}}

% The originals embedded the solution into an `ebox` and used \exrule
% as a horizontal divider between the exercise statement and the
% solution. Our new format moves the solution into a separate
% `.solution` div, so the divider should now collapse to nothing —
% defining \exrule as a no-op keeps any stray reference compilable.
\providecommand{\exrule}{\par}

% xcolor brings the dvipsnames palette (BrickRed, BurntOrange,
% ForestGreen, CadetBlue, RoyalPurple) the originals reach for.
\PassOptionsToPackage{dvipsnames,svgnames,x11names}{xcolor}

% MAPDA uses \overbar{…} for binary negation. Define it as a thin
% over-line wrapper (cleaner than \overline for boolean variables).
\providecommand{\overbar}[1]{\mkern 1.5mu\overline{\mkern-1.5mu#1\mkern-1.5mu}\mkern 1.5mu}
\usepackage[framemethod=TikZ]{mdframed}
\usepackage[most]{tcolorbox}
\tcbuselibrary{breakable, skins}
\usepackage{etoolbox}
\usepackage{fancyvrb}
\usepackage{fvextra}   % adds the breaklines/breakanywhere keys to \fvset
\usepackage{enumitem}
\setlist{itemsep=2pt, topsep=4pt, parsep=0pt}

% --- colour palette (matches the website's UniPD accent) --------------
\definecolor{unipdRed} {HTML}{9b0014}
\definecolor{ink}      {HTML}{1a1a1a}
\definecolor{inkSoft}  {HTML}{3d3d3d}
\definecolor{inkFaint} {HTML}{6b6b6b}
\definecolor{rule}     {HTML}{c8c8c8}
\definecolor{surface}  {HTML}{f5f5f5}

\color{ink}

% --- KOMA heading typography (slim, restrained) -----------------------
\addtokomafont{disposition}{\normalfont}
\setkomafont{section}      {\normalfont\Large\bfseries}
\setkomafont{subsection}   {\normalfont\large\bfseries}
\setkomafont{subsubsection}{\normalfont\normalsize\bfseries}
\setkomafont{paragraph}    {\normalfont\normalsize\itshape}

% --- mdframed styles (the only "decoration") --------------------------
% Left rule, no box. Two tints: accent (theorems) and ink (definitions).
\mdfdefinestyle{leftRuleAccent}{
  linewidth=1.2pt, linecolor=unipdRed,
  topline=false, bottomline=false, rightline=false,
  innertopmargin=2pt, innerbottommargin=4pt,
  innerleftmargin=12pt, innerrightmargin=0pt,
  skipabove=8pt, skipbelow=8pt,
  backgroundcolor=white,
}
\mdfdefinestyle{leftRuleInk}{
  style=leftRuleAccent, linecolor=ink,
}
\mdfdefinestyle{leftRuleFaint}{
  style=leftRuleAccent, linecolor=rule, linewidth=0.9pt,
}

% --- theorem environments --------------------------------------------
% The companion Lua filter (Notes/pdf/code-output.lua) converts every
% Div with one of these classes (.theorem, .proposition, .lemma,
% .corollary, .definition, .example, .exercise, .solution, .remark,
% .note) into raw  \begin{<env>}[<optional name>] ... \end{<env>}
% blocks AND strips the class so Pandoc no longer auto-defines anything.
% So this preamble owns the env definitions outright.

\makeatletter
\@ifundefined{c@chapter}{%
  \newcommand{\notesNumberWithin}{section}%
}{%
  \newcommand{\notesNumberWithin}{chapter}%
}
\makeatother

% amsthm styles: plain (italic body, bold head), definition (upright body,
% bold head), remark (upright body, italic head) — the standard look.

\theoremstyle{plain}
\newtheorem{theorem}{Theorem}[\notesNumberWithin]
\newtheorem{proposition}[theorem]{Proposition}
\newtheorem{lemma}[theorem]{Lemma}
\newtheorem{corollary}[theorem]{Corollary}

\theoremstyle{definition}
\newtheorem{definition}{Definition}[\notesNumberWithin]
% example, exercise, solution are tcolorbox-based — defined further below.
% `solution` and `remark` are auto-defined by Pandoc's LaTeX template
% (always emitted, regardless of usage); we override `solution` via
% \AtBeginDocument so our tcolorbox version wins.
% `note` is not auto-defined — keep ours.
\theoremstyle{remark}
\newtheorem*{note}{Note}

% ── tcolorbox-based exercise / example / solution ────────────────────
% Common look: small coloured title bar attached at top-left ("Exercise
% N.M" or "Exercise N.M (optional name)"), thin matching border, white
% body, breakable across pages. Different colour tones distinguish them:
%   * exercise → accent (UniPD red)   — for problems to solve
%   * example  → ink (near-black)     — for worked illustrations
%   * solution → faint (light grey)   — for solutions to exercises

\newcommand{\notesTcbTitle}[3]{%
  % #1 = label ("Exercise"), #2 = counter text, #3 = optional name
  #1~#2\if\relax\detokenize{#3}\relax\else\ \normalfont\itshape(#3)\fi%
}

\newtcolorbox[auto counter, number within=\notesNumberWithin]{exercise}[1][]{%
  enhanced, breakable,
  colback=white,
  colframe=unipdRed,
  boxrule=0.5pt, arc=2pt, outer arc=2pt,
  fonttitle=\bfseries, coltitle=white,
  attach boxed title to top left={xshift=10pt, yshift=-8pt},
  boxed title style={
    colback=unipdRed, colframe=unipdRed,
    boxrule=0pt, arc=2pt,
    boxsep=3pt, left=8pt, right=8pt,
  },
  title={\notesTcbTitle{Exercise}{\thetcbcounter}{#1}},
  top=12pt, bottom=8pt, left=12pt, right=10pt,
  before skip=10pt, after skip=10pt,
}

\newtcolorbox[auto counter, number within=\notesNumberWithin]{example}[1][]{%
  enhanced, breakable,
  colback=white,
  colframe=ink,
  boxrule=0.5pt, arc=2pt, outer arc=2pt,
  fonttitle=\bfseries, coltitle=white,
  attach boxed title to top left={xshift=10pt, yshift=-8pt},
  boxed title style={
    colback=ink, colframe=ink,
    boxrule=0pt, arc=2pt,
    boxsep=3pt, left=8pt, right=8pt,
  },
  title={\notesTcbTitle{Example}{\thetcbcounter}{#1}},
  top=12pt, bottom=8pt, left=12pt, right=10pt,
  before skip=10pt, after skip=10pt,
}

% Pandoc auto-emits `\newtheorem*{solution}`. We replace it after the
% fact with our tcolorbox version so it carries the same visual identity.
\AtBeginDocument{%
  \let\solution\relax\let\endsolution\relax%
  \newtcolorbox{solution}[1][]{%
    enhanced, breakable,
    colback=surface,
    colframe=rule,
    boxrule=0.4pt, arc=2pt, outer arc=2pt,
    fonttitle=\bfseries, coltitle=ink,
    attach boxed title to top left={xshift=10pt, yshift=-8pt},
    boxed title style={
      colback=surface, colframe=rule,
      boxrule=0.4pt, arc=2pt,
      boxsep=3pt, left=8pt, right=8pt,
    },
    title={Solution\if\relax\detokenize{#1}\relax\else\ \normalfont\itshape(#1)\fi},
    top=10pt, bottom=6pt, left=10pt, right=10pt,
    before skip=8pt, after skip=10pt,
  }%
}

% Wrap each env with the left-rule mdframed style.
\newcommand{\notesWrapAccent}[1]{%
  \BeforeBeginEnvironment{#1}{\begin{mdframed}[style=leftRuleAccent]}%
  \AfterEndEnvironment{#1}{\end{mdframed}}}

\newcommand{\notesWrapInk}[1]{%
  \BeforeBeginEnvironment{#1}{\begin{mdframed}[style=leftRuleInk]}%
  \AfterEndEnvironment{#1}{\end{mdframed}}}

\newcommand{\notesWrapFaint}[1]{%
  \BeforeBeginEnvironment{#1}{\begin{mdframed}[style=leftRuleFaint]}%
  \AfterEndEnvironment{#1}{\end{mdframed}}}

\notesWrapAccent{theorem}
\notesWrapAccent{proposition}
\notesWrapAccent{lemma}
\notesWrapAccent{corollary}

\notesWrapInk{definition}
% example, exercise and solution are tcolorbox-based — no mdframed wrap.

\notesWrapFaint{remark}
\notesWrapFaint{note}

% --- proof environment: italic header, square QED ---------------------
\renewcommand{\qedsymbol}{$\square$}

% --- code OUTPUT environment ------------------------------------------
% Used by the code-output.lua filter to wrap each `cell-output-*` Div.
% Visual: thin grey left rule, no box, 90% body size, slim mono font.
\definecolor{codeBg}{HTML}{f3f3f3}

\mdfdefinestyle{codeOutput}{
  linewidth=1.4pt, linecolor=inkFaint,
  topline=false, bottomline=false, rightline=false,
  innertopmargin=2pt, innerbottommargin=4pt,
  innerleftmargin=10pt, innerrightmargin=0pt,
  skipabove=2pt, skipbelow=8pt,
  backgroundcolor=codeBg,
}
\newenvironment{codeoutput}
  {\begin{mdframed}[style=codeOutput]\footnotesize\ttfamily\linespread{1.05}\selectfont}
  {\end{mdframed}}

% Auto-wrap long lines so code and output never overflow the right
% margin. fvextra is loaded above so the keys are defined now.
\fvset{
  breaklines=true,
  breakanywhere=true,
  breakautoindent=true,
  breaksymbolleft={\hbox{\scriptsize\textcolor{inkFaint}{\ensuremath{\hookrightarrow}}\,}},
  breaksymbolsepleft=2pt,
}

% --- math display: tighter matrix / vector spacing --------------------
% A future iteration will hook these into the lua filter for ASCII →
% LaTeX matrix conversion. For now we just nudge the defaults so manual
% \begin{pmatrix}…\end{pmatrix} blocks look tidier.
\setlength{\arraycolsep}{4pt}
\renewcommand{\arraystretch}{1.1}

% --- spacing & list defaults ------------------------------------------
\setlength{\parskip}{0.4em plus 0.1em minus 0.05em}
\setlength{\parindent}{0pt}

% Quarto handles hyperref + link colours from the YAML (colorlinks,
% linkcolor, urlcolor, citecolor). No extra \hypersetup needed here.

% --- baked-in watermark for downloadable PDFs -------------------------
% draftwatermark stamps a diagonal "site copy" string onto every page as
% part of the page content (not as a separate annotation), so removal
% requires re-rendering from source rather than a tooling pass. The
% NOTES_NO_WATERMARK environment variable lets the user temporarily
% disable it during local previews:  NOTES_NO_WATERMARK=1 quarto render
\IfFileExists{draftwatermark.sty}{%
  \usepackage{draftwatermark}%
  \SetWatermarkText{\textcopyright\ alvaromendezrt@gmail.com -- site copy}%
  \SetWatermarkScale{0.55}%
  \SetWatermarkLightness{0.88}%
  \SetWatermarkAngle{45}%
}{}
